\chapter{Trabalhos Relacionados}
\label{sec:trabalhos-relacionados}

Neste capítulo, são apresentadas pesquisas recentes que abordam o problema da recuperação multimodal em documentos complexos, destacando seus pontos de convergência com a presente proposta, bem como as limitações que motivam o desenvolvimento do EDGAR.

\section{ColPali: Efficient Document Retrieval with Vision Language Models}
\label{sec:colpali}

\citeonline{faysse2025colpali} propõem o \textit{ColPali}, uma arquitetura de recuperação de documentos que substitui os tradicionais \textit{pipelines} de extração textual por um modelo de linguagem e visão (VLM) capaz de indexar páginas inteiras como representações visuais. O sistema opera por meio de um mecanismo de interação tardia (\textit{late interaction}), que compara vetores associados a \textit{patches} visuais da página com o vetor da consulta do usuário. Com isso, a abordagem reduz a dependência de OCR e diminui o custo de indexação. Os resultados apresentados pelos autores indicam desempenho competitivo em \textit{benchmarks} de recuperação de documentos com \textit{layouts} complexos, incluindo tabelas e infográficos.

Em termos de convergência técnica, o EDGAR adota o paradigma de indexação baseado em Modelos de Linguagem e Visão (VLM), explorado pelo \textit{ColPali}, ao empregar representações visuais em sua camada macro para capturar a semântica do \textit{layout} sem depender exclusivamente da linearização do texto. Contudo, enquanto o \textit{ColPali} concentra seu fluxo operacional na recuperação de páginas completas, o que pode introduzir ruído visual periférico e reduzir a precisão do contexto fornecido ao modelo gerador \citeb{yun2025lilac}, o EDGAR amplia essa base ao associar a indexação macro a um \textit{pipeline} de granularidade fina. Por meio da análise de layout de documentos (DLA) e do isolamento geométrico de blocos, o presente trabalho segmenta e refina o contexto didático em unidades menores, como parágrafos, tabelas e figuras. Dessa forma, a arquitetura proposta busca atender à necessidade de maior especificidade no domínio educacional, preservando o elo pedagógico entre a explicação teórica e o elemento visual correspondente.

\section{LILaC: Late Interacting in Layered Component Graph for Open-domain Multimodal Multihop Retrieval}
\label{sec:lilac}

\citeonline{yun2025lilac} apresentam o LILaC, um \textit{framework} de recuperação multimodal que estrutura documentos como grafos de componentes em camadas, interligando explicitamente parágrafos, tabelas e imagens por meio de arestas semânticas e estruturais. Essa representação em grafo permite que o sistema realize consultas que exigem o cruzamento de informações distribuídas em múltiplas páginas, alcançando alta precisão em tarefas de recuperação multimodal multi-salto (\textit{multihop}) em domínio aberto.

A principal convergência entre o EDGAR e o LILaC está na fragmentação estrutural do documento em seus componentes constitutivos, mapeando de forma explícita a relação entre blocos textuais e elementos visuais para preservar a integridade do contexto original. Contudo, as arquiteturas diferem quanto à complexidade algorítmica e à estratégia de modelagem dessas relações. Enquanto o LILaC estabelece essa interligação por meio de grafos dinâmicos em camadas e requer a decomposição de consultas com apoio de LLMs, o que aumenta a carga computacional e a latência no fluxo síncrono \citeb{yun2025lilac}, o EDGAR propõe uma alternativa voltada à leveza do \textit{pipeline}. A arquitetura proposta busca preservar a hierarquia semântica e pedagógica do material didático por meio de uma indexação híbrida, vinculando os recortes geométricos diretamente aos metadados dos vetores no fluxo \textit{offline}. Essa decisão de design contribui para fornecer densidade contextual ao modelo gerador, favorecendo a mitigação de alucinações semânticas sem comprometer a baixa latência e a responsividade esperadas na interface do Luminia Reader. Além disso, o EDGAR especializa-se no domínio educacional e em suas particularidades estruturais, diferentemente do LILaC, que é orientado a cenários de recuperação multimodal em domínio aberto.

\section{SLEUTH: Resolving Evidence Sparsity via Agentic Context Engineering for Long-Document Understanding}
\label{sec:sleuth}

\citeonline{liu2025resolving} propõem o SLEUTH, um \textit{framework} baseado em engenharia de contexto agêntica para a compreensão de documentos longos. O sistema emprega múltiplos agentes de inteligência artificial colaborativos que refinam iterativamente a busca por evidências esparsas ao longo do documento, avaliando a complexidade das consultas, extraindo pistas textuais e visuais e consolidando os fragmentos recuperados em um contexto coerente para o modelo gerador. Essa abordagem demonstrou resultados expressivos em tarefas de compreensão que exigem raciocínio sobre evidências distribuídas em documentos extensos.

O presente trabalho aproxima-se do SLEUTH ao empregar a engenharia de contexto como estratégia para refinar as evidências fornecidas ao modelo gerador, priorizando a curadoria dos fragmentos recuperados em vez da simples inserção de dados brutos. Contudo, as soluções diferem quanto à temporalidade e à natureza computacional dessa curadoria. Enquanto o SLEUTH opera de forma reativa e dinâmica no fluxo síncrono (\textit{online}), delegando a filtragem contextual a uma infraestrutura iterativa baseada em múltiplos agentes autônomos, o que pode resultar em maior latência e custo computacional por consulta, o EDGAR trata esse problema como uma etapa de engenharia estrutural integrada ao \textit{pipeline} assíncrono (\textit{offline}). Ao refinar o contexto previamente por meio da análise de layout e do isolamento geométrico de blocos, o presente trabalho busca fornecer ao modelo gerador um fragmento informacional mais denso e focado. Essa decisão arquitetural reduz a necessidade de orquestrações agênticas em tempo de execução, favorecendo a baixa latência e a responsividade necessárias para a integração prática do ecossistema ao leitor \textit{web} Luminia Reader.

\section{Análise Comparativa}
\label{subsec:analise-comparativa}

A Tabela~\ref{tab:comparativo_trabalhos} sintetiza os principais eixos de diferenciação entre os trabalhos relacionados e a arquitetura proposta neste estudo. Os critérios de comparação foram selecionados com base nas dimensões mais relevantes para o domínio educacional: a precisão da recuperação, o custo computacional associado e a capacidade de lidar com a estrutura multimodal e pedagógica dos materiais didáticos em PDF.

\begin{table}[htbp]
\centering
\caption{Análise comparativa dos trabalhos relacionados e o presente trabalho.}
\label{tab:comparativo_trabalhos}
\resizebox{\textwidth}{!}{%
\begin{tabular}{lccccc}
\toprule
\textbf{Critério} & 
\textbf{ColPali} & 
\textbf{LILaC} & 
\textbf{SLEUTH} & 
\textbf{EDGAR} \\
\midrule
Granularidade de recuperação & 
Página inteira & 
Componente via grafo & 
Trecho por agente & 
Componente isolado \\

Custo computacional & 
Baixo & 
Alto & 
Alto & 
Baixo \\

Suporte a elementos visuais & 
Sim & 
Sim & 
Sim & 
Sim \\

Preserva relações entre componentes & 
Não & 
Sim & 
Parcial & 
Sim \\

Orquestração de múltiplos agentes & 
Não & 
Não & 
Sim & 
Não \\

Mitigação proativa de alucinações & 
Não & 
Não & 
Sim (Síncrona) & 
Sim (Assíncrona)\\

Foco em domínio educacional & 
Não & 
Não & 
Não & 
Sim \\

Requer OCR & 
Não & 
Sim & 
Não & 
Sim \\
\bottomrule
\end{tabular}%
}
\fonte{Elaborado pelo autor.}
\end{table}

A análise evidencia que, embora os trabalhos relacionados apresentem contribuições relevantes para a recuperação multimodal, eles não conciliam simultaneamente baixo custo computacional, granularidade fina de recuperação e sensibilidade à estrutura pedagógica dos documentos. O ColPali, apesar de sua eficiência, opera em granularidade de página inteira, o que pode introduzir ruído visual no contexto fornecido ao modelo. O LILaC e o SLEUTH alcançam maior precisão em seus respectivos cenários de avaliação, porém dependem de arquiteturas computacionalmente mais intensivas, cuja latência acumulada pode comprometer a interatividade em tempo real exigida por um protótipo de leitura web. O EDGAR propõe-se a ocupar essa lacuna ao combinar recuperação em nível de componente isolado, baixo custo computacional e curadoria estrutural do contexto, características orientadas às demandas de validação do sistema proposto.